\chapter{Putting It Together: The Theory of the Paper, Step by Step}\label{sec:L5}

The first four lectures built tools. This lecture assembles them into the
argument of the paper. We go in the order the paper does: the model
(\cref{L5:sec:model}), where the variance of prefill work comes from
(\cref{L5:sec:variance}), the decision problem and its single price
(\cref{L5:sec:decision}), and the scheduler whose every step compares that
price with the price of memory (\cref{L5:sec:scheduler}). We close with an
honest list of what is exact and what is approximate
(\cref{L5:sec:honest}). For each step we name the result of Lectures 1--4
that carries it.

\section{The model in one page}\label{L5:sec:model}

\begin{definition}[Replica and its three demands]\label{L5:def:replica}
The system has $J$ \term{replicas}; each has a memory pool of $M$ bytes.
A turn $i$ has context $K_i$ tokens (already in the session), appends $n_i$
new tokens and decodes $o_i$ output tokens. It makes three demands on the
replica that serves it.
\begin{enumerate}[label=(\alph*)]
\item \emph{Prefill compute.} Prefilling $n$ tokens on top of a resident
  prefix of $K$ tokens costs
  \begin{equation}\label{L5:eq:prefill}
    P(n,K)=a\,n+b\,n\,(K+n/2)
  \end{equation}
  seconds, with $a,b\ge0$: a linear term for the dense layers and an
  attention term over the prefix. A \term{hit} costs $P(n_i,K_i)$; a
  \term{miss} recomputes the whole context, $P(K_i+n_i,0)$.
\item \emph{Decode bandwidth.} $D_i=o_i(\beta K_i+\omega/m)$ seconds, where
  $m$ is the batch size, $\beta$ the time to read one token's KV and
  $\omega$ the time to read the weights once.
\item \emph{Memory.} $c_i=\kappa K_i$ bytes while the state is resident,
  with $\kappa$ the KV bytes per token.
\end{enumerate}
\end{definition}

The hit/miss distinction lives entirely in demand (a): $D_i$ and $c_i$ are
the same for a hit and a miss. By \cref{cor:price-mono} the extra work of a
miss is $\Delta S_i=aK_i+bK_i^2/2$, quadratic in the context.

\paragraph{Partial reuse in the executable model.}
With a reusable prefix of $C\in[0,K]$ tokens, the turn computes
$P(K+n-C,C)$ rather than choosing only the two endpoints.
The additional work over a full hit is $r_K(K-C)$ from
\cref{L4:sec:knapsack}. In a block pool $C$ is rounded down to full
reusable blocks and cannot cover the final prompt token used for logits.
For such runs report reusable tokens and the distribution of $C$;
a binary hit rate alone does not determine the first two work moments.
The binary variance identity remains true for its stated mixture, but its
conditional means must be recomputed when partial misses are grouped together.

\paragraph{Two stations and a pool.} \Cref{fig:projection} in
\cref{L1:sec:projection} shows the structure; here we add the numbers and
name the lecture that backs each modelling choice.
\begin{enumerate}
\item \emph{Decode is a PS station} with capacity $\varphi(m)$
  (\cref{def:ps}). Continuous batching advances every turn in the batch at
  each step. By \cref{thm:insens} the batch-size law is
  $\pi(m)\propto\rho_D^m/\prod_{k\le m}\varphi(k)$ with $\rho_D=\lambda\E[D]$,
  whatever the law of $D$.
\item \emph{Prefill is a FIFO queue} whose server is available the fraction
  of time the decode batch leaves. Chunked prefill interleaves chunks with
  decode steps but serves pending prefills in arrival order, so a
  Poisson, fixed-service approximation uses \cref{sec:L3}, with the PK wait
  $\E[W_q]=\lambda\E[S^2]/(2(1-\rho))$ (\cref{thm:pk}).
\item \emph{Sessions form a closed loop} (prefill $\to$ decode $\to$ tool),
  and new sessions arrive as a Poisson process of rate $\Lambda$. The tool
  call is a delay station (\cref{L4:def:delay}).
\item \emph{The pool} of $M$ bytes holds both the running batch and the
  paused sessions' KV. Once it is full, the KV policy decides what to
  drop, and so which future turns hit.
\end{enumerate}

\paragraph{The objective.} The paper minimises $L=L_P+L_D$, the mean number
of turns in the prefill stage plus the mean number in the decode batch. By
Little's law (\cref{thm:little}), $L$ equals $\lambda$ times the mean time a
turn spends in the two stages, so minimising $L$ is minimising mean delay.

\serving{The two stations respond to the cache in opposite ways. The decode
stage sees only the mean demand (insensitivity), and a miss does not change
the decode demand at all (\cref{L2:cor:decode-hit}). The prefill stage sees
the second moment of its work (PK), and the cache sets that second moment.
This asymmetry is the whole reason the paper prices misses at the prefill
queue.}

\section{Where the variance of prefill work comes from}\label{L5:sec:variance}

The PK formula makes $\E[S^2]$, equivalently $\mathrm{CV}^2$, a first-class
quantity (\cref{cor:variance}). The paper does not assume that agentic
workloads are intrinsically ``high variance''; it asks where the variance
comes from. The tool is the law of total variance.

\begin{lemma}[Variance split of the hit/miss mixture]\label{L5:lem:split}
Let $S_h$ be the mixture of \cref{L3:def:mixture}: $S^{\mathrm{hit}}$ with
probability $h$ and $S^{\mathrm{miss}}$ with probability $1-h$, with means
$\mu_{\mathrm{hit}},\mu_{\mathrm{miss}}$ and variances
$\sigma^2_{\mathrm{hit}},\sigma^2_{\mathrm{miss}}$. Then
\[
  \Var[S_h]=\underbrace{h\,\sigma^2_{\mathrm{hit}}+(1-h)\,\sigma^2_{\mathrm{miss}}}_{\text{spread of the appends}}
  +\underbrace{h(1-h)\,(\mu_{\mathrm{miss}}-\mu_{\mathrm{hit}})^2}_{\text{hit/miss mixture}} .
\]
\end{lemma}

\begin{proof}
Let $I$ be the indicator of a hit. By the law of total variance,
$\Var[S_h]=\E[\Var(S_h\mid I)]+\Var(\E[S_h\mid I])$. The first term is
$h\sigma^2_{\mathrm{hit}}+(1-h)\sigma^2_{\mathrm{miss}}$. The conditional
mean $\E[S_h\mid I]$ takes the value $\mu_{\mathrm{hit}}$ with probability
$h$ and $\mu_{\mathrm{miss}}$ with probability $1-h$, a two-point
variable whose variance is $h(1-h)(\mu_{\mathrm{miss}}-\mu_{\mathrm{hit}})^2$.
\end{proof}

The first term is not the scheduler's to change: it comes from the spread
of tool outputs (a line or a whole file). The second term is the KV
policy's lever. It vanishes at $h\in\{0,1\}$ and grows with the squared gap
$(\mu_{\mathrm{miss}}-\mu_{\mathrm{hit}})^2$, which by \eqref{L5:eq:prefill}
grows like $K^4$ once attention dominates. The paper measures both terms
on production traces rather than assuming either one dominates. The lesson
for the theory is that what a scheduler must price is not a ratio of
$\mathrm{CV}^2$ but the rise of $\lambda\E[S^2]/(2(1-\rho))$ that one
more miss causes. That is exactly $\Phi_i$.

\section{The decision problem and one price}\label{L5:sec:decision}

At a decision epoch a replica holds a set of paused programs. Program
$i$ holds $c_i$ bytes, resumes with probability $p_i$ after an expected
further $\tau_i$ seconds, and its next turn needs $S_i^{\mathrm{hit}}$ if
its state is still resident and $S_i^{\mathrm{hit}}+\Delta S_i$ if not. The
controller decides for each state whether to keep it, and the memory
constraint couples all these decisions.

\paragraph{Step 1: price a miss.} \Cref{thm:price} gives the rise in $L_P$
when a fraction of turns turns from hits into misses:
\[
  \Phi_i=\underbrace{\Delta S_i}_{\text{own work}}
  +\underbrace{\frac{\lambda\bigl((S_i^{\mathrm{miss}})^2-(S_i^{\mathrm{hit}})^2\bigr)}{2(1-\rho)}}_{\text{head-of-line}}
  +\underbrace{\frac{\lambda W\,\Delta S_i}{1-\rho}}_{\text{load}} .
\]
A single eviction is a transient event, while \cref{thm:price} compares two
stationary regimes. The paper therefore reads one eviction as a small
change of policy, a fraction $q_i$ of arrivals turning from hits into
misses. This is an approximation, and the paper tests it empirically.

\paragraph{Step 2: price a drop.} Dropping program $i$'s state causes a
miss only if the program resumes, so the expected price of the drop, its
price per byte and its price per byte-second held are
\begin{equation}\label{L5:eq:utility}
  w_i=p_i\,\Phi_i,\qquad v_i=\frac{w_i}{c_i},\qquad u_i=\frac{w_i}{c_i\,\tau_i}.
\end{equation}

\paragraph{Step 3: price memory.} Keeping program $i$ occupies $c_i$ bytes
for about $\tau_i$ seconds, that is $c_i\tau_i$ byte-seconds. Let
$x_i\in\{0,1\}$ indicate ``keep'' and let $B$ be the byte-seconds available
over the horizon. The keep-or-drop problem is
\begin{equation}\label{L5:eq:primal}
  \mathrm{OPT}=\min_{x\in\{0,1\}^I}\ \sum_i(1-x_i)\,w_i
  \quad\text{s.t.}\quad \sum_i c_i\tau_i\,x_i\le B .
\end{equation}
Relaxing the constraint with a multiplier $\theta\ge0$ gives the
Lagrangian
$\mathcal{L}(x,\theta)=\sum_i(1-x_i)w_i+\theta\bigl(\sum_ic_i\tau_ix_i-B\bigr)$.
We call $\theta$ the \term{shadow price} of a byte-second.

\begin{proposition}[One price for memory]\label{L5:prop:shadow}
\begin{enumerate}[label=(\roman*)]
\item For every $\theta\ge0$, $x^\theta_i=\1\{u_i>\theta\}$ minimises
  $\mathcal{L}(\cdot,\theta)$, and $g(\theta):=\mathcal{L}(x^\theta,\theta)\le\mathrm{OPT}$.
\item If $\theta$ is such that $x^\theta$ uses the budget exactly,
  $\sum_ic_i\tau_ix^\theta_i=B$, then $x^\theta$ is optimal
  for~\eqref{L5:eq:primal}.
\end{enumerate}
\end{proposition}

\begin{proof}
(i) $\mathcal{L}(x,\theta)=\sum_iw_i-\theta B+\sum_ix_i(\theta c_i\tau_i-w_i)$
is separable in $x$. Each term is minimised by $x_i=1$ when
$\theta c_i\tau_i<w_i$, that is $u_i>\theta$, and by $x_i=0$ otherwise.
For any feasible $x$ the bracket $\sum_ic_i\tau_ix_i-B$ is $\le0$, so
$\sum_i(1-x_i)w_i\ge\mathcal{L}(x,\theta)\ge g(\theta)$. Minimising over
feasible $x$ gives $\mathrm{OPT}\ge g(\theta)$ (weak duality).

(ii) $x^\theta$ is feasible, and its constraint term vanishes, so its
objective equals $g(\theta)\le\mathrm{OPT}$. Hence it is optimal.
\end{proof}

Part~(ii) is the threshold rule of \cref{thm:threshold} in byte-second form:
keep what is worth more per byte-second than memory, drop the rest. When
no $\theta$ uses the budget exactly (items are indivisible), the fractional
relaxation is still solved by the threshold rule with one item split, and
this is a mathematical relaxation. The tail-block cost described in
\cref{L4:sec:knapsack} does not satisfy its linear-cost assumption.

\paragraph{Step 4: how long to keep.} \Cref{L5:prop:shadow} decides at one instant.
Because the Lagrangian is separable, once $\theta$ is fixed each program can also be
treated on its own over time: keep its state until some time $T$ and drop it then if the
program has not resumed. Let the program resume with probability $p\in(0,1]$, and, given
that it resumes, after a time with a continuous distribution function $G$ and density
$g$ on $[0,\infty)$. Holding the state costs $\theta c$ per second; dropping it costs
$\Phi$ if the program resumes later. The expected cost of the rule ``hold until $T$'' is
\begin{equation}\label{L5:eq:hold-cost}
  C(T)=\theta c\int_0^T\bigl(1-pG(t)\bigr)\dd t+\Phi\,p\,\bigl(1-G(T)\bigr),
\end{equation}
since the state is held at time $t$ exactly when the program has not resumed by $t$,
and a miss occurs exactly when it resumes after $T$. The quantity that decides is the
\term{resume hazard}
\[
  \eta(t)=\frac{p\,g(t)}{1-pG(t)},
\]
the rate at which a program that has not resumed by $t$ resumes at $t$.

\begin{proposition}[Optimal holding time]\label{L5:prop:hold}
Let $c>0$, $\theta>0$ and $r=\theta c/\Phi$ with $\Phi>0$.
Suppose $g$ is continuous. Define $\eta(T)=0$ after $pG(T)=1$.
\begin{enumerate}[label=(\roman*)]
\item $C'(T)=\bigl(1-pG(T)\bigr)\bigl(\theta c-\Phi\,\eta(T)\bigr)$ for every $T\ge0$.
  Holding a little longer lowers the cost exactly when $\Phi\,\eta(T)>\theta c$.
\item If $\eta$ is nonincreasing, let $T^*=\inf\{T\ge0:\eta(T)\le r\}$. If
  $T^*<\infty$ it minimises $C$ over $[0,\infty)$; if $T^*=\infty$, $C$ is nonincreasing
  and the state should be held until the program resumes. In particular $T^*=0$, drop at
  once, if and only if $\eta(0)\le r$.
\item If $G$ is exponential with mean $\tau$, then $\eta$ is nonincreasing and
  $\eta(0)=p/\tau$, so the state should be kept at all if and only if
  \[
    u=\frac{p\,\Phi}{c\,\tau}>\theta ,
  \]
  the price per byte-second of \eqref{L5:eq:utility}. If $p=1$ and $u>\theta$, the
  state is held until the program resumes. If $p<1$ and $u>\theta$, the optimal holding
  time is
  \[
    T^*=\tau\,\ln\Bigl[\frac{p}{1-p}\Bigl(\frac{\Phi}{\theta c\tau}-1\Bigr)\Bigr].
  \]
\end{enumerate}
\end{proposition}

\begin{proof}
(i) The integrand in \eqref{L5:eq:hold-cost} is continuous, so by the fundamental
theorem of calculus the first term has derivative $\theta c(1-pG(T))$, and the second
has derivative $-\Phi p\,g(T)$. Factoring out $1-pG(T)$ gives the formula when
$pG(T)<1$; when $pG(T)=1$ (possible only if $p=1$) both sides vanish, because
$g(T)=0$ wherever $G$ has reached $1$ and $g$ is continuous. The factor $1-pG(T)$ is
nonnegative, so the sign of $C'(T)$ is the sign of $\theta c-\Phi\eta(T)$.

(ii) For $T<T^*$, $\eta(T)>r$, so $C'(T)\le0$ by (i). For $T>T^*$, $\eta(T)\le r$
because $\eta$ is nonincreasing, so $C'(T)\ge0$. Hence $C$ is nonincreasing on
$[0,T^*]$ and nondecreasing on $[T^*,\infty)$, and $T^*$ is a minimiser. If
$T^*=\infty$, then $C'\le0$ everywhere. The last claim is the definition of $T^*$ at
$T=0$, using that $\eta$ is nonincreasing.

(iii) With $x=e^{-T/\tau}$, $g(T)=x/\tau$ and $1-pG(T)=1-p+px$, so
\[
  \eta(T)=\frac{p}{\tau}\cdot\frac{x}{1-p+px}.
\]
The map $x\mapsto x/(1-p+px)$ is nondecreasing on $(0,1]$, and $x$ decreases in $T$,
so $\eta$ is nonincreasing; it is the constant $1/\tau$ when $p=1$. At $T=0$,
$\eta(0)=p/\tau$, and $\eta(0)>r$ is $p\Phi/(c\tau)>\theta$. If $p=1$, $\eta\equiv1/\tau>r$,
so $T^*=\infty$ by (ii). If $p<1$, $\eta(T)\to0$ as $T\to\infty$, so $T^*$ is finite and
solves $\eta(T^*)=r$ by continuity:
$px/\tau=r(1-p+px)$, that is $x=r(1-p)/\bigl(p(1/\tau-r)\bigr)$, and
$T^*=-\tau\ln x=\tau\ln\bigl[\frac{p}{1-p}\bigl(\frac{1}{r\tau}-1\bigr)\bigr]$ with
$1/(r\tau)=\Phi/(\theta c\tau)$.
\end{proof}

When $\theta=0$, holding until resumption has zero cost and is optimal.
For positive memory price, the instantaneous test $u_i>\theta$ of the scheduler is exactly the optimal decision
to keep a state for exponential resume times, and part~(iii) adds how long to keep it.
For example, with $p=0.5$, $\tau=10$\,s, $\Phi=20$\,s, $c=1$ and $\theta=0.3$, $u=1>\theta$,
the optimal holding time is $T^*\approx17.3$\,s, and it lowers the expected cost from
$C(0)=p\Phi=10$ to $C(T^*)\approx5.6$. When the gap distribution has a decreasing hazard, as
heavy-tailed gaps typically do, part~(ii) gives a program-specific time-to-live.

\begin{remark}[Unknown resume times and the oracle]\label{L5:rem:ski}
If nothing is known about $G$, holding the state is the classical \emph{ski rental}
problem: pay rent $\theta c$ per second or buy once at price $\Phi$
\cite{karlin1988}. The rule ``hold for $\Phi/(\theta c)$ seconds'' costs at most twice
the cost in hindsight for any resume time: if the program resumes at
$\tau\le\Phi/(\theta c)$ it pays $\theta c\tau$, which is optimal; otherwise it pays
$2\Phi$ against a hindsight optimum of $\Phi$. The hindsight optimum, which knows
whether and when each program resumes, is called the \term{oracle}; it keeps a state
exactly when the program resumes and $\theta c\tau<\Phi$. It cannot be run online, but
it is the benchmark against which a policy's cost is reported as cost/OPT. The factor
two needs the program to resume: against a program that never does, the oracle pays
nothing, and no rule that holds the state for a positive time has a bounded ratio.
\end{remark}

\Cref{L5:prop:hold} holds $\theta$ and $\Phi$ fixed. In the real system both move with
the policy: what is dropped changes the load, and the load changes every $\Phi_i$ and
the value of $\theta$ that fills the pool. The joint optimum is a Markov decision problem
over the whole state, with no closed form; \cref{L5:prop:hold} is its solution at a
fixed point of $(\theta,\Phi)$.

\section{The scheduler, step by step}\label{L5:sec:scheduler}

Every decision now compares a program's $u_i$ with the replica's $\theta$.

\paragraph{Eviction.} When $\Delta C$ bytes must be freed now, the problem
is the covering knapsack of \cref{L4:def:knapsack} with $w_i=p_i\Phi_i$.
For additive whole-program items, order the candidates by the corresponding
price and apply the covering-knapsack rule, with its indivisibility caveat.
The byte-second key $u_i$ instead prices retention over time. Block tail
costs are concave (\cref{L4:sec:knapsack}), so neither ordering alone proves
optimality for a block eviction. Recompute a candidate tail's lost value.
Deployed rules evict by a fixed key: shortest context first, least recently
used, longest idle, or an expired time-to-live. \Cref{prop:blind} says what
such a key can and cannot do: a key that ignores $p_i$ has no constant
approximation ratio (part~(i)), and shortest-first is within a factor two
only when every program resumes and the price is the recompute work alone
(part~(ii)).

\begin{example}[The load reorders programs]\label{L5:ex:flip}
Take $a=0.2$\,ms/token and $b=a/K_c$ with $K_c=50{,}000$ tokens, and an
append of $n=1000$ tokens. Program A has $K=200$k tokens and $p=0.3$;
program B has $K=20$k tokens and $p=0.9$. Then $\Delta S_A=120$\,s and
$\Delta S_B=4.8$\,s. At an idle queue ($\lambda\to0$) the price is the own
work, and the price per thousand tokens $p\Phi/K$ is $0.18$\,s for A and
$0.216$\,s for B: evict A first. At $\lambda=0.25$/s with $\E[S]=2$\,s and
$\E[S^2]=20$\,s$^2$ ($\rho=0.5$, $W=5$\,s) the head-of-line term dominates;
$\Phi_A\approx4080$\,s and $\Phi_B\approx23.2$\,s, and the prices per
thousand tokens become $6.12$\,s and $1.05$\,s: evict B first. No fixed key
reproduces both orders, which is why the scheduler prices at decision time.
\end{example}

\paragraph{Keep, offload or drop.} With a slower tier (host memory or a
shared KV store) each paused state has three options: keep it at cost
$\theta c_i\tau_i$, drop it at expected cost $w_i$, or offload it at the
\term{priced fetch}, which is $\Phi_i$ with the transfer time in place of
$\Delta S_i$ (plus the tier's own wait once it saturates). The scheduler
takes the least of the three. Two facts follow from this being a minimum.

\begin{proposition}[Offloading as an option]\label{L5:prop:offload}
Let each state $i$ have a finite set of actions $\mathcal{A}_i$ with
costs $\gamma_i(a)$, and let the total cost be separable. Then
\begin{enumerate}[label=(\roman*)]
\item enlarging each $\mathcal{A}_i$ never raises the optimal cost; in
  particular, making offloading available never hurts;
\item a policy that always offloads can cost more than one that never
  offloads.
\end{enumerate}
\end{proposition}

\begin{proof}
(i) A minimum over a larger set is at most the minimum over a subset.
(ii) One state with keep cost $1$, drop cost $1$ and priced fetch $2$
suffices: always offloading pays $2$, never offloading pays $1$.
\end{proof}

So when an always-offloading stack collapses, the collapse is a property of
the policy, not of the offloading mechanism.

\paragraph{Placement.} A resuming turn can run where its state is (replica
1) or elsewhere (replica $j$). The router charges each replica its wait $W_j$
and service $S_j$, the cost $R_j$ of moving or recomputing the state, and
the expected future cost $F_j$ of the program's remaining turns, and picks
the minimum. The price of a miss supplies $R_j$: moving turns the next turn
into a miss served at replica $j$, so $R_j$ is $\Phi_i$ at replica $j$'s
load. \emph{Session affinity}, the rule that never moves a program, sets
$R_j=\infty$ for $j\ne1$.

\begin{proposition}[When affinity loses]\label{L5:prop:rhostar}
Let the affinity replica be an M/M/1 queue with service rate $\mu$ and
utilisation $\rho<1$, and let the alternative be an idle replica with the
same rate, at extra cost $x=R+F\ge0$. Staying costs more than moving if and
only if
\[
  \rho>\rho^*:=\frac{\mu x}{1+\mu x},
\]
and $\rho^*$ is nondecreasing in $x$.
\end{proposition}

\begin{proof}
By \cref{prop:mm1} the mean response time at the affinity replica is
$1/(\mu(1-\rho))$, and moving costs $1/\mu+x$. The inequality
$1/(\mu(1-\rho))>1/\mu+x$ is $1/(1-\rho)>1+\mu x$, that is
$1-\rho<1/(1+\mu x)$, that is $\rho>\mu x/(1+\mu x)$. The map
$x\mapsto\mu x/(1+\mu x)=1-1/(1+\mu x)$ is nondecreasing.
\end{proof}

An idle alternative gives the earliest inversion, so $\rho^*$ is a lower
bound on the load at which affinity loses to any real alternative. A shared
KV store lowers $R$ and hence $\rho^*$: affinity and a store are
alternatives, not complements.

\paragraph{Admission.} The pool holds both the batch and the paused
states, and the two compete unstably: a miss keeps its turn longer in the
prefill queue and in the batch, which holds more bytes, which evicts more
paused state, which causes more misses. Admission keeps the replica out
of this regime by capping the number of live sessions. The cap comes from
\cref{cor:resident-kv}: with sessions of lifetime $T_s$ holding $K(t)$
tokens at age $t$, the mean resident KV is $\kappa\Lambda\E[\int_0^{T_s}K(t)\dd t]$.
Writing $\bar k=\kappa\E[\int_0^{T_s}K(t)\dd t]/\E[T_s]$ for the
time-averaged footprint of a live session and using Little's law for the
mean number of live sessions $\Lambda\E[T_s]$, the mean resident KV is
(live sessions)$\times\bar k$. The admission step caps live sessions near
$M/\bar k$ and holds the rest in an entry queue. \Cref{cor:resident-kv}
also warns that the naive footprint $\kappa\E[\bar K]$ underestimates
$\bar k$ whenever long contexts belong to long sessions. The cap has a price
of its own, the wait at the entry queue, and a session is worth admitting
while the memory it will hold, priced at $\theta$ per byte-second, costs
less than that wait.

\paragraph{The scheduler in one table.} Each deployed rule replaces one
price by a constant:
\begin{center}\small
\begin{tabular}{@{}ll@{}}
\toprule
Rule & What it fixes \\
\midrule
shortest-first eviction & $w_i\propto c_i^2$ (every program resumes, own work only) \\
LRU, time-to-live & $p_i$ constant \\
always offload & priced fetch $=0$ \\
strict session affinity & $R_j=\infty$ \\
admit every session & $\theta=0$ \\
\bottomrule
\end{tabular}
\end{center}

\section{What is exact and what is approximate}\label{L5:sec:honest}

A mathematician should know which statements are theorems about the model
and which are modelling choices.

\begin{itemize}
\item \emph{Exact, given the model:} the price bracket (\cref{thm:price}),
  the decode price and its monotonicity (\cref{thm:L-mono},
  \cref{prop:decode-price}), the finite-population inequalities
  (\cref{prop:finite}), the price-blind bounds (\cref{prop:blind}), and
  \cref{L5:prop:shadow,L5:prop:offload,L5:prop:rhostar}.
\item \emph{Approximate:} the prefill stage is not M/G/1. Session feedback
  makes its arrivals non-Poisson, and the server's budget fluctuates with
  the decode batch. BCMP product form covers the PS and delay stations but
  not the FIFO prefill stage. For exponential work,
  \cref{prop:finite}(ii) shows the open formula overstates the wait at
  equal utilisation. For general work this comparison is an empirical
  question.
\item \emph{Approximate:} one eviction is a transient event priced by a
  stationary comparison (Step~1 above).
\item \emph{Not in the queueing formula:} a turn whose KV does not fit
  waits for memory blocks before its prefill starts. That wait is set by
  the pool, not by the queue, and adds to the PK wait. The capacity
  function $\varphi$ depends on the footprint law and the admission rule
  and must be calibrated, not derived.
\end{itemize}

\paragraph{What the testbed showed about the last point.} On the paper's testbed
with contexts of about 50k tokens, the wait for the first token was a wait for memory
blocks: the prefill server was busy a small fraction of the time, and the PK wait was
one to two orders of magnitude below the observed one. Three engine rules made the
pool a queue. A waiting turn is admitted only when its whole prompt fits; admission is
in arrival order, so a turn that does not fit blocks everyone behind it (head-of-line
blocking, now on memory); and every admission takes the least recently used cached
blocks. What was observed: the turns that missed had the same think times as the hits,
and nearly all of them had arrived while their rank already had a queue. What a
post hoc model of the replica suggests, but the measurements do not show directly: the
admissions ahead of them evicted their prefixes while they waited. Either way the
misses were selected by congestion, so on that testbed the ratio of miss to hit waiting
times measures congestion, not the price of a miss. In the short-context runs the extra
misses caused by forced ones followed long think gaps and a higher rate of re-inserted
contexts, which points to the pool (re-insertions shortening how long a cached prefix
survives) rather than to the wait.
In queueing terms, the hit rate is not a policy input but depends on the wait, and the
wait on the hit rate; the feedback model in \cref{L6:sec:feedback} treats this dependence.

The paper tests the approximate links in two ways: in simulation on
replayed production sessions, and on a real testbed. The criterion is
\emph{decision faithfulness}: the model must rank policies the way the
real system does, even where its latencies are off.

\section{Exercises}

\begin{exercise}\label{L5:ex:split}
In \cref{L5:lem:split} let $\sigma_{\mathrm{hit}}=\sigma_{\mathrm{miss}}=\sigma$.
Find the hit rate $h$ that maximises the mixture term, and show that the
mixture term exceeds the append term exactly when
$h(1-h)(\mu_{\mathrm{miss}}-\mu_{\mathrm{hit}})^2>\sigma^2$.
\emph{Hint:} $h(1-h)$ is maximised at $h=1/2$.
\end{exercise}

\begin{exercise}\label{L5:ex:rhostar}
With $\mu=1$ turn/s, compute $\rho^*$ of \cref{L5:prop:rhostar} for
$x=0.1$, $1$ and $10$\,s, and interpret: how loaded must the affinity
replica be before moving a program with a ten-second recompute pays off?
\emph{Hint:} $\rho^*=x/(1+x)$ when $\mu=1$.
\end{exercise}
